\documentclass[aspectratio=169]{beamer}
\usepackage[T1]{fontenc}
\usepackage[utf8]{inputenc}
\usepackage{lmodern}
\usepackage[ngerman]{babel}
\usepackage{texstage}
\usepackage{texstage-transactions-schedule}
\setbeamertemplate{navigation symbols}{}
\setbeamertemplate{footline}[frame number]
\begin{document}
\begin{frame}{Transaktionsabläufe: Verzahnungen verstehen}
\begin{enumerate}
\item Endwerte vorhersagen.
\item Gespeicherte Lesewerte und Schreibschritte verfolgen.
\item Direkte Konflikte mit Zeugen erklären.
\item Den Ablauf unter Erhalt der Programmreihenfolge reparieren.
\end{enumerate}
\end{frame}
\begin{frame}
\begin{transactionsschedule}[id=counter]{.98\linewidth}{.90\textheight}
\end{transactionsschedule}
\end{frame}
\begin{frame}
\begin{transactionsschedule}[id=independent]{.98\linewidth}{.90\textheight}
\end{transactionsschedule}
\end{frame}
\begin{frame}
\begin{transactionsschedule}[id=read_read]{.98\linewidth}{.90\textheight}
\end{transactionsschedule}
\end{frame}
\begin{frame}
\begin{transactionsschedule}[id=write_write]{.98\linewidth}{.90\textheight}
\end{transactionsschedule}
\end{frame}
\begin{frame}
\begin{transactionsschedule}[id=chain]{.98\linewidth}{.90\textheight}
\end{transactionsschedule}
\end{frame}
\begin{frame}
\begin{transactionsschedule}[id=blind_write]{.98\linewidth}{.90\textheight}
\end{transactionsschedule}
\end{frame}
\end{document}
